\chapter{三维医学图像分割相关理论与关键技术}

% TODO(P-002): Cover only theory that later chapters actually use.

\section{三维医学图像与分割任务}

第一章已从研究背景、技术进展与现有不足出发，明确了本文围绕三维医学图像分割开展后续研究。本章进一步转向方法章节所依赖的基本理论与技术概念，本节首先限定后续网络实际处理的数据对象及其预测任务。这里不再讨论医学图像分割的应用意义或不同方法的优劣，而是从三维影像的数据组织方式出发，明确切片、体素、体数据以及体素级分割之间的关系，为后续三维特征建模建立统一的术语和符号基础。

CT、MRI等医学影像在计算处理中可表示为按空间顺序组织的二维切片（slice）或由连续切片构成的三维体数据（volume）。单张切片由二维空间位置上的像素组成；当切片按照深度方向排列后，三维空间中的基本离散单元通常称为体素（voxel）。因此，切片可以理解为三维体数据在某一方向上的二维截面，而体素则对应体数据中的一个三维空间位置。每个体素既具有空间坐标，也携带该位置上的影像强度或其他成像信号；分割标签则进一步为相同空间位置赋予离散的语义类别。一个体数据可用深度、高度和宽度三个空间维度描述，其中深度维反映切片沿第三个空间方向的排列关系，高度和宽度描述切片平面内的位置关系。已有三维医学分割研究直接面向 volumetric image 进行密集预测，3D U-Net将原二维U-Net中的相关操作扩展为三维对应形式，V-Net则以三维MRI体数据作为输入进行整幅体数据分割，说明三维体数据已经构成医学图像分割中的基本建模对象之一。\cite{cicek2016threeDunet}\cite{milletari2016vnet}

为了在后续章节中统一描述网络输入，可以把三维影像、参考标签和网络输出纳入同一张量表示。设输入影像为 \(\mathbf{X}\)，其通道数为 \(C\)，空间尺寸分别为 \(D\)、\(H\) 和 \(W\)；对应标签记为 \(\mathbf{Y}\)，共有 \(K\) 个前景类别并另设背景类。三维分割任务可统一表示为
\[
\begin{aligned}
&\mathbf{X}\in\mathbb{R}^{C\times D\times H\times W},\qquad
\mathbf{Y}\in\{0,1,\ldots,K\}^{D\times H\times W},\\
&f_{\theta}:\mathbf{X}\rightarrow\mathbf{P},\qquad
\mathbf{P}\in[0,1]^{(K+1)\times D\times H\times W},\qquad
\sum_{k=0}^{K}P_{k,d,h,w}=1,\\
&\hat{Y}_{d,h,w}=\arg\max_{k\in\{0,\ldots,K\}}P_{k,d,h,w}.
\end{aligned}
\]
其中，\(C\) 是通道维而不是空间维，\(D\)、\(H\) 和 \(W\) 分别表示深度、高度和宽度；对于单通道影像，\(C\) 可为1，当输入包含多个成像通道或模态时，也可以沿通道维组织相应信息。标签中的0表示背景，\(1,\ldots,K\) 表示待分割的目标类别。\(f_{\theta}\) 表示参数为 \(\theta\) 的分割网络，\(\mathbf{P}\) 表示各体素对应不同类别的预测概率，\(\hat{Y}_{d,h,w}\) 则为坐标 \((d,h,w)\) 处的最终预测标签。这里的符号只固定数据对象及输入输出关系，不预设本文后续实验采用的具体图像尺寸、体素间距、裁剪大小、网络实现或预处理方式。

由上述定义可以看出，三维医学图像分割属于空间密集预测任务。它与图像级分类的基本区别在于，网络并非只为整个输入体数据给出一个类别，而是需要保持输入与输出在三维空间位置上的对应关系，并为大量体素分别给出类别判断。对于二分类分割，通常只区分背景和一个目标类别，此时可令 \(K=1\)；对于多类别分割，则需要在同一标签图中区分背景与多个目标类别。无论类别数量如何变化，最终结果都可以表示为与输入三维空间相对应的标签体数据。3D U-Net与V-Net均以三维体数据上的密集分割为基本任务形式，前者强调对 volumetric image 生成 dense 3D segmentation，后者直接学习完整MRI体数据的分割预测，为上述体素级任务描述提供了对应的研究实例。\cite{cicek2016threeDunet}\cite{milletari2016vnet}

三维任务的关键并不在于简单增加一个数组维度，而在于深度、高度和宽度共同确定了体素的空间位置。对于位于 \((d,h,w)\) 的体素，其空间关系既可以发生在同一切片内部，也可以发生在不同深度位置之间；连续切片因而不是一组在空间上彼此无关的二维图像，而是三维体数据的不同截面。将它们按空间顺序组织后，深度方向与平面内方向能够在统一坐标中共同描述，使跨切片的形态延续、位置变化等信息具备被联合建模的条件。\cite{cicek2016threeDunet}\cite{milletari2016vnet}这里所强调的是三维数据表示提供了利用这些空间关系的可能，并不意味着只要采用三维输入就能自动获得充分的上下文，也不意味着二维方法完全无法通过相邻切片、序列输入等其他方式引入额外信息。对于后续章节而言，“三维卷积”“长距离上下文”等概念所作用的对象，正是这种具有三维空间坐标和体素对应关系的数据或特征表示。

在具体计算过程中，三维分割既可以面向完整体数据，也可以处理从完整体数据中截取的局部三维体块。若完整体数据记为 \(\mathbf{X}\)，则局部三维 patch 或 crop 可以理解为从其空间范围中取得的子张量 \(\mathbf{X}_{p}\in\mathbb{R}^{C\times d\times h\times w}\)，其中 \(d\leq D\)、\(h\leq H\)、\(w\leq W\)。局部体块覆盖的空间范围虽然小于完整体数据，但仍保留深度、高度和宽度三个有序空间维度，其内部体素仍具有三维邻接与位置关系，因此相应预测仍属于三维体素级分割。V-Net展示了对完整MRI体数据进行整体预测的研究形式，而DeepMedic等三维CNN方法采用局部三维区域并结合多尺度输入完成分割，说明 whole-volume 与 patch-based 处理主要体现输入覆盖范围和计算组织方式的差异，而不会改变三维分割以空间体素为基本预测对象的任务性质。\cite{milletari2016vnet}\cite{kamnitsas2017deepmedic} 本节仅作概念区分，本文所采用的具体体数据处理方式将在后续实验设置中进一步说明。

综上，三维医学图像分割可以归纳为：以由体素组成的三维医学影像或其局部体块为输入，在深度、高度和宽度共同构成的空间中，为各体素预测对应的背景或目标类别，并输出与输入空间位置相对应的标签体数据。slice、voxel、volume、输入张量与标签图由此形成后续方法章节共享的数据语义接口。在明确网络面对的三维数据对象与体素级预测目标之后，下一步需要进一步说明模型如何在该三维空间中形成局部特征表示，并如何组织不同空间分辨率层级的信息，因此下一节将介绍卷积神经网络与U形编码器--解码器的相关基础。

\section{卷积神经网络与U形编码器--解码器}

\subsection{三维卷积与局部特征提取}

% TODO: Explain only concepts used in Work1 and Work2.

\subsection{U-Net编码器--解码器与跳跃连接}

% TODO: Connect multi-scale feature hierarchy to the Work2 problem setting.

\section{长距离上下文建模}

\subsection{Transformer与自注意力机制}

% TODO: Approved literature only.

\subsection{xLSTM与Vision-xLSTM}

% TODO: Use original sources; do not write formulas from memory.

\section{多尺度特征与尺度结构信息}

% TODO: Explain the problem background only; do not introduce SSE, IM, or SDE here.

\section{损失函数与评价指标}

\subsection{分割损失函数}

% TODO: Include only losses actually used by the final methods.

\subsection{分割精度指标}

% TODO: Dice and HD95; add others only when formally used.

\subsection{模型效率指标}

% TODO: Params, FLOPs, Peak VRAM, and inference time; distinguish their meanings.

\section{本章小结}

% TODO: Summarize the theory debt repaid for Chapters 3 and 4.
